\documentclass[aspectratio=169,10pt]{beamer}

% Shared Beamer setup for Fall 2026 course slides.
% Clean Cal Poly-inspired theme: no custom class files or uncommon packages.
\mode<presentation>{
  \usetheme{default}
  \usefonttheme{professionalfonts}
}

\definecolor{CalPolyGreen}{HTML}{154734}
\definecolor{CalPolyGold}{HTML}{BD8B13}
\definecolor{CalPolySage}{HTML}{7FA28F}
\definecolor{CalPolyMint}{HTML}{DDF1E7}
\definecolor{CalPolySand}{HTML}{A29F86}
\definecolor{CalPolyBeige}{HTML}{EFEEDE}
\definecolor{DeckInk}{HTML}{1F2933}
\definecolor{DeckMuted}{HTML}{5F6B7A}
\definecolor{DeckSoft}{HTML}{F7F7F5}

\setbeamersize{text margin left=0.55in,text margin right=0.55in}
\setbeamercolor{normal text}{fg=DeckInk,bg=white}
\setbeamercolor{structure}{fg=CalPolyGreen}
\setbeamercolor{title}{fg=white}
\setbeamercolor{frametitle}{fg=CalPolyGreen,bg=white}
\setbeamercolor{block title}{bg=CalPolySage,fg=white}
\setbeamercolor{block body}{bg=CalPolyMint,fg=black}
\setbeamercolor{alerted text}{fg=CalPolyGold}

\setbeamerfont{title}{size=\Huge,series=\bfseries}
\setbeamerfont{frametitle}{size=\Large,series=\bfseries}
\setbeamerfont{block title}{size=\normalsize,series=\bfseries}

\setbeamertemplate{navigation symbols}{}
\setbeamertemplate{itemize item}{\color{CalPolyGold}$\blacktriangleright$}
\setbeamertemplate{itemize subitem}{\color{CalPolyGreen}$\bullet$}
\setbeamertemplate{footline}{
  \hfill{\color{DeckMuted}\scriptsize\insertframenumber/\inserttotalframenumber}
  \hspace{0.18in}\vspace{0.12in}
}
\setbeamertemplate{frametitle}{
  \vspace{0.18in}
  \hspace{0.02in}\insertframetitle\par
  \vspace{0.08in}
  \noindent\makebox[\linewidth][l]{%
    {\color{CalPolyGold}\rule{0.55in}{2pt}}%
    {\color{CalPolyGreen}\rule{\dimexpr\linewidth-0.55in\relax}{2pt}}%
  }
}

\newcommand{\CourseNumber}{Course}
\newcommand{\CourseTitle}{Course Title}
\newcommand{\LectureNumber}{0}
\newcommand{\LectureTitle}{Lecture Title}
\newcommand{\LectureDate}{Date}
\newcommand{\CourseUnit}{Unit}

\newcommand{\coursenumber}[1]{\renewcommand{\CourseNumber}{#1}}
\newcommand{\coursetitle}[1]{\renewcommand{\CourseTitle}{#1}}
\newcommand{\lecturenumber}[1]{\renewcommand{\LectureNumber}{#1}}
\newcommand{\lecturetitle}[1]{\renewcommand{\LectureTitle}{#1}}
\newcommand{\lecturedate}[1]{\renewcommand{\LectureDate}{#1}}
\newcommand{\courseunit}[1]{\renewcommand{\CourseUnit}{#1}}

\author{Kirk Duran}
\institute{California Polytechnic State University, San Luis Obispo}
\date{\LectureDate}

\newcommand{\makedecktitle}{
  \begingroup
  \setbeamercolor{background canvas}{bg=CalPolyGreen}
  \begin{frame}[plain]
    \vfill
    \begin{center}
      {\usebeamerfont{title}\color{white}\LectureTitle\par}
      \vspace{0.18in}
      {\Large\color{white}Lecture \LectureNumber\par}
    \end{center}
    \vfill
    \noindent{\color{white}\rule{\linewidth}{1.2pt}}\par
    \vspace{0.08in}
    \noindent\colorbox{CalPolyGold}{\makebox[0.24\linewidth][c]{\color{white}\Large\CourseNumber}}
    \hspace{0.12in}{\color{white}\Large Kirk Duran}
  \end{frame}
  \endgroup
}

\newenvironment{keyidea}{\begin{block}{Key idea}}{\end{block}}
\newenvironment{activity}{\begin{block}{In class}}{\end{block}}
\newenvironment{cpdefinition}{\begin{block}{Definition}}{\end{block}}
\newenvironment{exampleblockcp}{
  \begingroup
  \setbeamercolor{block title}{bg=CalPolySand,fg=white}
  \setbeamercolor{block body}{bg=CalPolyBeige,fg=black}
  \begin{block}{Example}
}{
  \end{block}
  \endgroup
}

\newcommand{\imageplaceholder}[2][0.3in]{%
  \noindent\makebox[\linewidth][c]{%
    \fcolorbox{CalPolySage}{CalPolyBeige}{%
      \parbox[c][#1][c]{0.86\linewidth}{%
        \centering
        {\color{DeckMuted}\scriptsize Image placeholder: }%
        {\color{DeckInk}\scriptsize\ttfamily\detokenize{#2}}%
      }%
    }%
  }%
}

\newcommand{\sectionframe}[1]{
  \begingroup
  \setbeamercolor{background canvas}{bg=CalPolyGreen}
  \begin{frame}[plain]
    \vfill
    {\color{white}\LARGE\bfseries #1\par}
    \vspace{0.18in}
    {\color{CalPolyGold}\rule{0.22\paperwidth}{3pt}}
    {\color{white}\rule{0.62\paperwidth}{3pt}}
    \vfill
  \end{frame}
  \endgroup
}

\newcommand{\placeholdernote}{
  \begin{block}{Placeholder}
    Replace these bullets with the final lecture material, examples, and in-class activities.
  \end{block}
}


\pdfstringdefDisableCommands{\def\translate#1{#1}}

\graphicspath{{../assets/lecture-19/}}

% If an image file is not yet available, the slide displays a labeled
% placeholder using the intended filename. Place the matching file in
% ../assets/lecture-19/ to replace the placeholder automatically.
\newcommand{\lectureimage}[2][1.75in]{%
  \IfFileExists{../assets/lecture-19/#2}{%
    \includegraphics[width=\linewidth,height=#1,keepaspectratio]{#2}%
  }{%
    \fbox{%
      \parbox[c][#1][c]{0.94\linewidth}{%
        \centering
        \small\textbf{Image Placeholder}\\[0.08in]
        \scriptsize\texttt{\detokenize{#2}}
      }%
    }%
  }%
}

\setbeamersize{text margin left=0.42in,text margin right=0.42in}

\coursenumber{CS 1001}
\coursetitle{Introduction to Computing}
\lecturenumber{19}
\lecturetitle{Strings as Sequences}
\lecturedate{Fall 2026}
\courseunit{7. Sequential Data}

\begin{document}

\makedecktitle

\begin{frame}{Today}
  \begin{itemize}
    \item Model a string as an ordered sequence of characters.
    \item Create strings using quotation marks and escape sequences.
    \item Access characters using positive and negative indexes.
    \item Measure and slice strings using the same positional model used for lists.
    \item Explain string immutability and distinguish mutation from rebinding.
    \item Combine, compare, search, split, join, and transform strings.
    \item Connect Python strings to Unicode code points.
  \end{itemize}

  \begin{keyidea}
    A string is an immutable sequence: its character positions can be read and sliced, but an existing string value cannot be changed in place.
  \end{keyidea}
\end{frame}

\sectionframe{Part 1: Strings Are Sequential Data}

\begin{frame}{Definition: String}
  \begin{cpdefinition}
    A \textbf{string} is a Python sequence value representing text. A string contains an ordered sequence of Unicode characters.
  \end{cpdefinition}

  \begin{block}{Example}
    \centering
    \Large\texttt{text = "Python"}
  \end{block}

  \begin{itemize}
    \item The sequence has a fixed order.
    \item Each position contains one character.
    \item Character positions are identified by integer indexes.
    \item The empty string contains zero characters: \texttt{""}.
    \item Text such as \texttt{"42"} is still a string; Python has no separate built-in character type.
  \end{itemize}
\end{frame}

\begin{frame}{The Same Sequence Model Applies to Lists and Strings}
  \begin{columns}[T,onlytextwidth]
    \begin{column}{0.47\textwidth}
      \begin{block}{List}
        \centering
        \texttt{[10, 20, 30, 40]}
      \end{block}
      \begin{itemize}
        \item ordered positions
        \item integer indexes
        \item slicing
        \item mutable
      \end{itemize}
    \end{column}
    \begin{column}{0.47\textwidth}
      \begin{block}{String}
        \centering
        \texttt{"code"}
      \end{block}
      \begin{itemize}
        \item ordered positions
        \item integer indexes
        \item slicing
        \item immutable
      \end{itemize}
    \end{column}
  \end{columns}

  \begin{keyidea}
    Lists and strings share positional sequence operations, but differ fundamentally in whether the existing value may be changed.
  \end{keyidea}
\end{frame}

\begin{frame}{A Useful Mental Model: Adjacent Character Positions}
  \begin{columns}[T,onlytextwidth]
    \begin{column}{0.52\textwidth}
      \begin{block}{String}
        \centering
        \Large\texttt{"Python"}
      \end{block}
      \begin{center}
        \renewcommand{\arraystretch}{1.3}
        \begin{tabular}{c|c|c|c|c|c|c}
          index & 0 & 1 & 2 & 3 & 4 & 5 \\\hline
          char. & P & y & t & h & o & n
        \end{tabular}
      \end{center}
    \end{column}
    \begin{column}{0.42\textwidth}
      \lectureimage[2.20in]{string-adjacent-positions.png}
    \end{column}
  \end{columns}

  \begin{keyidea}
    For tracing, think of a string as consecutive character positions numbered from left to right.
  \end{keyidea}
\end{frame}

\begin{frame}{Creating String Literals}
  \begin{block}{Single-line strings}
    \texttt{name = "Ada"}\\
    \texttt{language = 'Python'}
  \end{block}

  \begin{block}{Multi-line literal}
    \texttt{message = """Line one}\\
    \texttt{Line two"""}
  \end{block}

  \begin{itemize}
    \item Single and double quotation marks both delimit ordinary string literals.
    \item Matching quotation marks must close the literal.
    \item Triple quotation marks may span physical source-code lines.
  \end{itemize}
\end{frame}


\sectionframe{Part 2: Indexing and Length}

\begin{frame}{String Indexing Uses Zero-Based Positions}
  \begin{block}{Example}
    \texttt{text = "Python"}\\[0.08in]
    \texttt{text[0]} $\longrightarrow$ \texttt{"P"}\\
    \texttt{text[1]} $\longrightarrow$ \texttt{"y"}\\
    \texttt{text[5]} $\longrightarrow$ \texttt{"n"}
  \end{block}

  \begin{keyidea}
    Indexing a string produces a new one-character string containing the character at that position.
  \end{keyidea}
\end{frame}

\begin{frame}{Negative Indexes Count from the End}
  \begin{block}{Example}
    \texttt{text = "Python"}
  \end{block}

  \begin{center}
    \renewcommand{\arraystretch}{1.3}
    \begin{tabular}{c|c|c|c|c|c|c}
      character & P & y & t & h & o & n \\\hline
      positive & 0 & 1 & 2 & 3 & 4 & 5 \\
      negative & -6 & -5 & -4 & -3 & -2 & -1
    \end{tabular}
  \end{center}

  \begin{block}{Examples}
    \texttt{text[-1]} $\longrightarrow$ \texttt{"n"}\qquad
    \texttt{text[-2]} $\longrightarrow$ \texttt{"o"}
  \end{block}
\end{frame}

\begin{frame}{An Invalid String Index Raises \texttt{IndexError}}
  \begin{block}{Example}
    \texttt{text = "cat"}
  \end{block}

  \begin{itemize}
    \item Valid positive indexes are \texttt{0}, \texttt{1}, and \texttt{2}.
    \item Valid negative indexes are \texttt{-1}, \texttt{-2}, and \texttt{-3}.
    \item \texttt{text[3]} raises \texttt{IndexError}.
    \item \texttt{text[-4]} also raises \texttt{IndexError}.
  \end{itemize}

  \begin{keyidea}
    Indexing requires one existing position. The same boundary rule applies to lists and strings.
  \end{keyidea}
\end{frame}

\begin{frame}{\texttt{len()} Counts Characters}
  \begin{cpdefinition}
    \texttt{len(string)} returns the number of character positions in the string.
  \end{cpdefinition}

  \begin{block}{Examples}
    \texttt{len("Python")} $\longrightarrow$ \texttt{6}\\
    \texttt{len("")} $\longrightarrow$ \texttt{0}\\
    \texttt{len("Hello, world!")} $\longrightarrow$ \texttt{13}
  \end{block}

  \begin{itemize}
    \item Spaces and punctuation count as characters.
    \item Length is not the final positive index.
  \end{itemize}
\end{frame}

\begin{frame}{Length and the Final Index}
  \begin{block}{Example}
    \texttt{text = "Python"}\\[0.08in]
    \texttt{len(text)} $\longrightarrow$ \texttt{6}\\
    \texttt{text[len(text) - 1]} $\longrightarrow$ \texttt{"n"}
  \end{block}

  \begin{keyidea}
    For a nonempty string, the final positive index is \texttt{len(text) - 1}.
  \end{keyidea}
\end{frame}


\sectionframe{Part 3: String Slicing}

\begin{frame}{String Slicing Uses the Same \texttt{start:stop:step} Model}
  \begin{block}{General syntax}
    \centering
    \Large\texttt{text[start:stop:step]}
  \end{block}

  \begin{itemize}
    \item \texttt{start} is included.
    \item \texttt{stop} is excluded.
    \item \texttt{step} controls the distance between selected positions.
    \item The result is a new string.
  \end{itemize}
\end{frame}

\begin{frame}{The Stop Boundary Is Still Exclusive}
  \begin{block}{Example}
    \texttt{text = "Python Programming"}\\[0.10in]
    \texttt{text[0:6]} $\longrightarrow$ \texttt{"Python"}
  \end{block}

  \begin{itemize}
    \item Index \texttt{0} is selected.
    \item Indexes \texttt{1} through \texttt{5} are selected.
    \item Index \texttt{6} is the stopping boundary and is not selected.
  \end{itemize}

  \begin{keyidea}
    Read \texttt{0:6} as ``start at 0, stop before 6.''
  \end{keyidea}
\end{frame}

\begin{frame}{Omitting Slice Bounds Uses Sequence Boundaries}
  \begin{block}{Examples}
    \texttt{text = "Python Programming"}\\[0.08in]
    \texttt{text[:6]} $\longrightarrow$ \texttt{"Python"}\\
    \texttt{text[7:]} $\longrightarrow$ \texttt{"Programming"}\\
    \texttt{text[:]} $\longrightarrow$ \texttt{"Python Programming"}
  \end{block}

  \begin{itemize}
    \item The original string remains unchanged.
    \item A slice expression evaluates to another string value.
  \end{itemize}
\end{frame}

\begin{frame}{The Optional Step Changes Which Positions Are Selected}
  \begin{block}{Example}
    \texttt{text = "abcdefgh"}\\[0.08in]
    \texttt{text[0:8:2]} $\longrightarrow$ \texttt{"aceg"}\\
    \texttt{text[1:8:2]} $\longrightarrow$ \texttt{"bdfh"}
  \end{block}

  \begin{itemize}
    \item The default step is \texttt{1}.
    \item A step of \texttt{2} selects every second position within the slice bounds.
    \item A step of \texttt{0} is invalid and raises \texttt{ValueError}.
  \end{itemize}
\end{frame}

\begin{frame}{A Negative Step Produces Characters in Reverse Position Order}
  \begin{block}{Example}
    \texttt{text = "Python"}\\[0.08in]
    \texttt{text[::-1]} $\longrightarrow$ \texttt{"nohtyP"}
  \end{block}

  \begin{columns}[T,onlytextwidth]
    \begin{column}{0.53\textwidth}
      \begin{itemize}
        \item Omitted bounds adapt to the negative direction.
        \item The result is a new reversed string.
        \item The original \texttt{text} remains \texttt{"Python"}.
      \end{itemize}
    \end{column}
    \begin{column}{0.40\textwidth}
      \lectureimage[1.65in]{string-reverse-slice.png}
    \end{column}
  \end{columns}
\end{frame}

\begin{frame}{Practice: Predict the Slice Results}
  \begin{block}{Given}
    \centering
    \texttt{text = "abcdefghij"}
  \end{block}

  \begin{activity}
    Predict each result before evaluating it.
    \begin{enumerate}
      \item \texttt{text[2:6]}
      \item \texttt{text[:4]}
      \item \texttt{text[6:]}
      \item \texttt{text[1:9:3]}
      \item \texttt{text[-4:-1]}
      \item \texttt{text[::-1]}
    \end{enumerate}
  \end{activity}
\end{frame}

\sectionframe{Part 4: Immutability and Building New Strings}

\begin{frame}{Strings Are Immutable}
  \begin{cpdefinition}
    An \textbf{immutable} value cannot be changed after that value has been created.
  \end{cpdefinition}

  \begin{block}{Invalid operation}
    \texttt{text = "cat"}\\
    \texttt{text[0] = "b"}
  \end{block}

  \begin{alertblock}{Result}
    Python raises \texttt{TypeError}: individual string positions cannot be assigned new characters.
  \end{alertblock}
\end{frame}

\begin{frame}{Concatenation Creates a New String}
  \begin{block}{Example}
    \texttt{first = "Hello"}\\
    \texttt{second = " world"}\\
    \texttt{message = first + second}
  \end{block}

  \begin{block}{Values afterward}
    \texttt{first} $\longrightarrow$ \texttt{"Hello"}\\
    \texttt{second} $\longrightarrow$ \texttt{" world"}\\
    \texttt{message} $\longrightarrow$ \texttt{"Hello world"}
  \end{block}

  \begin{keyidea}
    The \texttt{+} operator combines string values into a new string; it does not modify either operand.
  \end{keyidea}
\end{frame}

\begin{frame}{Concatenation Requires String Operands}
  \begin{block}{Invalid}
    \texttt{"Age: " + 21}
  \end{block}

  \begin{block}{Explicit conversion}
    \texttt{"Age: " + str(21)} $\longrightarrow$ \texttt{"Age: 21"}
  \end{block}

  \begin{itemize}
    \item \texttt{str(value)} produces a string representation of a value.
    \item Converting to a string changes the type of the produced value, not the original variable.
  \end{itemize}
\end{frame}

\begin{frame}{Reassignment Is Not Mutation}
  \begin{block}{Code}
    \texttt{text = "Python"}\\
    \texttt{text = "C" + text[1:]}
  \end{block}

  \begin{itemize}
    \item \texttt{text[1:]} creates the new string \texttt{"ython"}.
    \item Concatenation creates the new string \texttt{"Cython"}.
    \item Assignment then rebinds the name \texttt{text} to that new value.
    \item The original string value \texttt{"Python"} was never altered.
  \end{itemize}

  \begin{keyidea}
    A variable can be rebound even when the value it previously referred to is immutable.
  \end{keyidea}
\end{frame}


\begin{frame}{Escape Sequences Represent Special Characters Inside a Literal}
  \begin{block}{Common examples}
    \texttt{"line one\textbackslash nline two"} \quad newline\\
    \texttt{"name\textbackslash tvalue"} \quad tab\\
    \texttt{"She said \textbackslash"hello\textbackslash"."} \quad embedded quotation marks\\
    \texttt{"C:\textbackslash\textbackslash Users\textbackslash\textbackslash Ada"} \quad literal backslashes
  \end{block}

  \begin{itemize}
    \item An escape sequence is written using a backslash in source code.
    \item The resulting string contains the represented character, not necessarily the literal source characters used to write the escape.
  \end{itemize}
\end{frame}

\sectionframe{Part 5: Membership, Equality, and Character Codes}

\begin{frame}{Membership Checks Whether a Substring Occurs}
  \begin{block}{Examples}
    \texttt{text = "Python programming"}\\[0.08in]
    \texttt{"Python" in text} $\longrightarrow$ \texttt{True}\\
    \texttt{"gram" in text} $\longrightarrow$ \texttt{True}\\
    \texttt{"python" in text} $\longrightarrow$ \texttt{False}\\
    \texttt{"Java" not in text} $\longrightarrow$ \texttt{True}
  \end{block}

  \begin{itemize}
    \item Membership can search for one character or an entire substring.
    \item The comparison is case-sensitive.
    \item Both \texttt{in} and \texttt{not in} produce Boolean values.
  \end{itemize}
\end{frame}


\begin{frame}{String Equality Requires the Same Characters in the Same Order}
  \begin{block}{Examples}
    \texttt{"Hello" == "Hello"} $\longrightarrow$ \texttt{True}\\
    \texttt{"Hello" == "hello"} $\longrightarrow$ \texttt{False}\\
    \texttt{"cat" == "cat "} $\longrightarrow$ \texttt{False}\\
    \texttt{"abc" != "acb"} $\longrightarrow$ \texttt{True}
  \end{block}

  \begin{keyidea}
    Equality compares string contents exactly; case, order, spaces, and punctuation matter.
  \end{keyidea}
\end{frame}

\begin{frame}{Ordering Comparisons Are Lexicographic}
  \begin{cpdefinition}
    String ordering compares corresponding characters from left to right until a difference determines the result; if one string is a prefix of the other, the shorter string orders first.
  \end{cpdefinition}

  \begin{block}{Examples}
    \texttt{"apple" < "banana"} $\longrightarrow$ \texttt{True}\\
    \texttt{"app" < "apple"} $\longrightarrow$ \texttt{True}
  \end{block}

  \begin{itemize}
    \item Character ordering is based on Unicode code points.
    \item Therefore, ordering is case-sensitive and may not match dictionary conventions used by humans.
  \end{itemize}
\end{frame}

\begin{frame}{Python Strings Use Unicode}
  \begin{itemize}
    \item Python's \texttt{str} type represents Unicode text.
    \item Unicode assigns each character a numerical \textbf{code point}.
    \item ASCII characters occupy the first 128 Unicode code points.
  \end{itemize}

  \begin{block}{Examples}
    \texttt{ord("A")} $\longrightarrow$ \texttt{65}\\
    \texttt{ord("Z")} $\longrightarrow$ \texttt{90}\\
    \texttt{ord("a")} $\longrightarrow$ \texttt{97}\\
    \texttt{chr(65)} $\longrightarrow$ \texttt{"A"}
  \end{block}

  \begin{keyidea}
    Because uppercase and lowercase letters have different code points, relational comparisons are case-sensitive; for example, \texttt{"Zebra" < "apple"} is \texttt{True}.
  \end{keyidea}
\end{frame}


\sectionframe{Part 6: Converting and Transforming Strings}

\begin{frame}{\texttt{list(string)} Produces a List of One-Character Strings}
  \begin{block}{Example}
    \texttt{text = "hello"}\\
    \texttt{chars = list(text)}\\[0.08in]
    \texttt{chars} $\longrightarrow$ \texttt{["h", "e", "l", "l", "o"]}
  \end{block}

  \begin{itemize}
    \item The original string remains unchanged.
    \item The result is a new list.
    \item Because lists are mutable, the resulting character positions can later be changed as list elements.
  \end{itemize}
\end{frame}

\begin{frame}{\texttt{split()} Produces a List of Substrings}
  \begin{block}{Default whitespace separation}
    \texttt{text = "hello world Python"}\\
    \texttt{text.split()} $\longrightarrow$ \texttt{["hello", "world", "Python"]}
  \end{block}

  \begin{block}{Explicit separator}
    \texttt{"red,green,blue".split(",")}\\
    $\longrightarrow$ \texttt{["red", "green", "blue"]}
  \end{block}

  \begin{keyidea}
    \texttt{split()} does not modify the source string; it returns a new list containing string pieces.
  \end{keyidea}
\end{frame}

\begin{frame}{\texttt{join()} Combines a List of Strings}
  \begin{block}{Example}
    \texttt{words = ["hello", "world", "Python"]}\\
    \texttt{sentence = " ".join(words)}\\[0.08in]
    \texttt{sentence} $\longrightarrow$ \texttt{"hello world Python"}
  \end{block}

  \begin{itemize}
    \item The string before \texttt{.join(...)} is the separator.
    \item The list supplied to \texttt{join()} must contain strings.
    \item The result is a new string.
  \end{itemize}
\end{frame}

\begin{frame}{Case-Conversion Methods Return New Strings}
  \begin{block}{Examples}
    \texttt{text = "hello world"}\\[0.08in]
    \texttt{text.upper()} $\longrightarrow$ \texttt{"HELLO WORLD"}\\
    \texttt{text.lower()} $\longrightarrow$ \texttt{"hello world"}\\
    \texttt{text.capitalize()} $\longrightarrow$ \texttt{"Hello world"}\\
    \texttt{text.title()} $\longrightarrow$ \texttt{"Hello World"}
  \end{block}

  \begin{alertblock}{Important}
    None of these methods mutates \texttt{text}; each produces a new string value.
  \end{alertblock}
\end{frame}

\begin{frame}{Formatted String Literals Build Strings from Expressions}
  \begin{block}{Example}
    \texttt{name = "Ada"}\\
    \texttt{age = 36}\\
    \texttt{message = f"Name: \{name\}, Age: \{age\}"}
  \end{block}

  \begin{block}{Result}
    \texttt{message} $\longrightarrow$ \texttt{"Name: Ada, Age: 36"}
  \end{block}

  \begin{itemize}
    \item Expressions inside braces are evaluated first.
    \item Their formatted representations are inserted into the new string.
    \item Format specifications may control representation, e.g. \texttt{f"\{3.14159:.2f\}"} produces \texttt{"3.14"}.
    \item Formatting changes the produced text, not the underlying numeric value.
  \end{itemize}
\end{frame}


\begin{frame}{Integrated Trace: New Strings, No Mutation}
  \begin{block}{Code}
    \texttt{text = "python"}\\
    \texttt{head = text[:2]}\\
    \texttt{tail = text[2:]}\\
    \texttt{label = head.upper() + tail}\\
    \texttt{words = label.split("h")}\\
    \texttt{joined = "-".join(words)}
  \end{block}

  \begin{activity}
    Determine the final values of \texttt{text}, \texttt{head}, \texttt{tail}, \texttt{label}, \texttt{words}, and \texttt{joined}. Identify every expression that creates a new value.
  \end{activity}
\end{frame}

\begin{frame}{Strings in Memory: Conceptual Model vs. Implementation}
  \begin{columns}[T,onlytextwidth]
    \begin{column}{0.54\textwidth}
      \begin{itemize}
        \item Our tracing model treats a string as adjacent character positions.
        \item Python guarantees the sequence behavior: indexing, slicing, length, equality, and immutability.
        \item The exact physical representation of Unicode strings is an implementation detail.
        \item An implementation may use compact internal encodings rather than one fixed-size memory cell per visible character.
      \end{itemize}
    \end{column}
    \begin{column}{0.40\textwidth}
      \lectureimage[2.35in]{string-conceptual-vs-implementation.png}
    \end{column}
  \end{columns}

  \begin{keyidea}
    Use the positional sequence model for reasoning about Python code; do not assume a specific byte layout unless the representation itself is the topic.
  \end{keyidea}
\end{frame}

\begin{frame}{Operation Summary}
  \begin{center}
    \renewcommand{\arraystretch}{1.25}
    \begin{tabular}{l|l}
      \textbf{Operation} & \textbf{Result} \\\hline
      \texttt{text[i]} & one-character string \\
      \texttt{text[start:stop:step]} & new string \\
      \texttt{len(text)} & integer length \\
      \texttt{a + b} & new concatenated string \\
      \texttt{sub in text} & Boolean \\
      \texttt{text.split(...)} & new list of strings \\
      \texttt{sep.join(parts)} & new string \\
      \texttt{text.upper()} / \texttt{lower()} & new string
    \end{tabular}
  \end{center}
\end{frame}

\begin{frame}{Takeaways}
  \begin{itemize}
    \item A string is an ordered, indexed, immutable sequence of Unicode characters.
    \item String indexing and slicing use the same positional rules as lists, including negative indexes and an exclusive stop boundary.
    \item Concatenation, slicing, case conversion, formatting, and related operations produce new string values.
    \item Reassigning a variable to a new string is not mutation of the old string.
    \item Membership, equality, and relational comparisons are case-sensitive.
    \item \texttt{split()} converts text into a list of string pieces; \texttt{join()} combines string pieces into a new string.
  \end{itemize}

  \begin{keyidea}
    When tracing string code, track values and bindings carefully: the string value itself never changes in place.
  \end{keyidea}
\end{frame}

\end{document}
